Multimodal Large Language Models (MLLMs) show strong generalization in Vision-and-Language Navigation, yet their 2D-centric visual encoders struggle with 3D spatial reasoning. Training-time feature distillation from pretrained geometry foundation models injects 3D priors without inference overhead, but the design conditions for effective transfer remain unexamined.

This thesis evaluates three design factors of geometry feature distillation on VLN-CE: (1) teacher selection across feed-forward and recurrent models, (2) information flow directionality comparing bidirectional and causal teachers, and (3) temporal stride of navigation history. Results reveal three findings. First, under standard final-layer distillation, Pi3, which decodes geometry directly from its final layers, substantially outperforms the non-distilled baseline, whereas DPT-based models yield marginal gains or negative transfer, highlighting alignment between decoding heads and distillation targets. Second, bidirectional and causal variants exhibit no consistent advantage, indicating that attention directionality is not a primary bottleneck. Third, wider temporal strides improve overall performance through longer history, while distillation gains narrow at wider strides due to reduced visual overlap. Overall, target representation alignment and temporal context outweigh attention directionality, motivating multi-layer distillation for DPT-based teachers and decoupled stride designs.
